\documentclass[11pt]{article}
\usepackage[T1]{fontenc}
\usepackage{amsmath,amssymb,amsthm}
\usepackage[margin=1.12in]{geometry}
\usepackage[colorlinks=true,allcolors=blue,
  pdftitle={Frobenius norm bounds for Newton polygons with CM field action},
  pdfauthor={Roy van Rijn}]{hyperref}
\newtheorem{theorem}{Theorem}
\newtheorem{corollary}[theorem]{Corollary}
\newtheorem{proposition}[theorem]{Proposition}
\theoremstyle{definition}
\newtheorem{example}[theorem]{Example}
\DeclareMathOperator{\Tr}{Tr}
\DeclareMathOperator{\Nm}{N}
\newcommand{\Q}{\mathbb Q}
\newcommand{\Z}{\mathbb Z}
\newcommand{\fp}{f_p}
\title{Frobenius norm bounds for Newton polygons\ProtectBreak with CM field action}
\newcommand{\ProtectBreak}{\texorpdfstring{\\}{ }}
\author{Roy van Rijn}
\date{}
\begin{document}
\maketitle

\begin{abstract}
Let an abelian variety over a number field have a Rosati-compatible
action of a CM field of degree \(2e\), and let \(h\) be the relative
rank of its first cohomology. If at least one conjugate Hodge pair is
mixed, we prove simultaneous strict bounds on the rational-prime
valuations of the norms of all intermediate relative Frobenius
coefficients. The bounds hold at all good absolute degree-one primes
outside a finite set after one finite base extension, and they constrain
every corresponding Newton polygon. In relative rank two with \(b>0\)
balanced pairs they give \(p\)-rank at least
\(\dim A-2\lfloor(b-1)/2\rfloor\). In relative rank four they give
density-one ordinary reduction after finite extension if exactly one
conjugate pair is mixed and its type is \((2,2)\). This holds in
arbitrarily large dimensions and with extra geometric endomorphisms.
A least polygon with prescribed zero and unit multiplicities gives
a closed formula for the sharp numerical minimum \(p\)-rank in every
signature. We classify precisely when it forces ordinarity and
realize the first surviving rank-two polygon over prime fields with
any prescribed auxiliary torsion level. In relative rank two with
three balanced pairs, linked projective monodromy gives density-one
ordinarity after finite extension. We also give infinite cyclic-cover
moment exclusions, including cases certified by a selected Frobenius
line on a nonordinary Jacobian factor.
\end{abstract}

\section{Introduction}

The ordinary-reduction conjecture attributed to Serre predicts a
positive-density set of ordinary primes for an abelian variety over a
number field \cite{fite}. A stronger form asks for density one after
finite extension \cite[\S7]{pink}\cite[\S1.1]{cantoral}. An action of
a CM field can make the ordinary condition accessible because at split
primes it separates Frobenius into smaller components. The geometric
endomorphism algebra can contain additional endomorphisms throughout
this paper.

There are two bounds on a relative Frobenius coefficient. Its complex
absolute values satisfy a Weil bound, while the slopes of each local
component give lower bounds on its valuations. When the latter make a
normalized field norm an integer, adjoining enough torsion forces it
to attain the upper Weil bound. Equality in the triangle inequality
then makes the eigenvalues within each relative component coincide.
A mixed split component cannot have a single repeated eigenvalue.
The trace would put that eigenvalue in \(\Q_p\), where its strictly
fractional slope is impossible. The same finite extension works for
each coefficient.

Fix an abelian variety \(A/K\) and an embedding
\(E\hookrightarrow\operatorname{End}^0(A_{\overline K})\) of a CM
field of degree \(2e\). Assume Rosati for some polarization restricts
to complex conjugation on \(E\). Set \(g=\dim A\) and \(h=g/e\).
Fix a complex embedding of the base field over which the action is
defined. Write \(d_\sigma\) for the multiplicity on the tangent space
at an embedding \(\sigma:E\hookrightarrow\mathbb C\). The conjugate
pair \((d_\sigma,d_{\bar\sigma})\) sums to \(h\), and is
\emph{mixed} when both entries are positive. Use the convention in
which the sum of crystalline slopes in a component is the dimension
of the corresponding \(p\)-divisible group. Interchanging all
embedding labels by conjugation has no effect on the statements.

\begin{samepage}
\begin{theorem}[Simultaneous coefficient and slope bounds]\label{thm:bounds}
Suppose there is a mixed pair. There is a finite extension \(L/K\) over
which the action and polarization are defined and which contains a
normal closure of \(E\). At every good absolute degree-one prime of
\(L\) outside a finite set, let \(p\) be the residue characteristic.
Write \(a_j\in E\) for the \(j\)-th elementary symmetric function of
relative Frobenius, and \(s_{\sigma,1}\le\cdots\le s_{\sigma,h}\) for
the slopes at the split places of \(E\) over \(p\). For every
\(1\le j<h\) simultaneously,
\[
 v_p\bigl(\Nm_{E/\Q}(a_j)\bigr)\le ej-1,\qquad
 \sum_\sigma\left\lceil\sum_{k=1}^j s_{\sigma,k}\right\rceil
 \le ej-1 .
\]
These primes have natural density one in \(L\).
\end{theorem}
\end{samepage}

\begin{corollary}[Rank two]\label{cor:ranktwo}
Suppose \(h=2\) and let \(b\) be the number of mixed pairs, which
then have type \((1,1)\). If \(b>0\), at the primes of
Theorem~\ref{thm:bounds} at most \(\lfloor(b-1)/2\rfloor\) of these
pairs are nonordinary. In particular
\[
 \fp(A)\ge g-2\left\lfloor\frac{b-1}{2}\right\rfloor .
\]
If \(b=0\), all split good primes outside the finite set are ordinary.
\end{corollary}

\begin{corollary}[Rank four]\label{cor:rankfour}
Suppose \(h=4\) and exactly one pair is mixed, of type \((2,2)\).
Ordinary primes have density one after finite extension. This holds
for every \(e\) and thus for dimension \(4e\).
\end{corollary}

\section{Relative coefficients and the norm argument}\label{sec:coefficients}

We first specify the local input and the comparison defining the
coefficients. Pass to a finite extension with the action and polarization
defined and containing a normal closure of \(E\).
The full order \(R=E\cap\operatorname{End}(A)\) has finite index in
\(\mathcal O_E\). Away from this index and finitely many excluded primes
it completes to the maximal order at the residue prime.
A degree-one prime of the enlarged base lies over a rational prime split
in \(E\), outside ramification. The idempotents of
\(\mathcal O_E\otimes\Z_p\) act on the \(p\)-divisible group of the
reduction. Their heights are \(h\), and their dimensions are the
characteristic-zero signature entries, since the Lie algebra of the
abelian scheme is locally free and decomposes by the same idempotents.
The sum of the slopes of a component is its dimension. The polarization
identifies the slopes of the conjugate component with \(1-s\) in reverse
order. These assertions can equivalently be read from the split
Dieudonn\'e modules and their polarizations.

At a good prime with residue field \(\mathbb F_p\) let \(\pi\) be its
\(p\)-power Frobenius. For an actual endomorphism \(u\) of the
reduction, the characteristic polynomial on \(H^1\) is the degree
polynomial of \(u\) in both prime-to-\(p\) and crystalline cohomology.
At any sufficiently large integer \(m\) its value is
\(\deg([m]-u)\), by the exterior algebra description of cohomology and
the action on top cohomology. Consequently for each positive integer
\(k\) the values \(\Tr(a\pi^k\mid H^1)\) agree in the two
cohomologies for every \(a\in R\).
The nondegenerate trace pairing on \(E/\Q\) identifies the relative
power sums with the scalar extensions of a single element of \(E\)
for each \(k\). The same comparison for \(a\) alone gives relative
rank \(h\), using the prime-to-\(p\) comparison with the Betti
module, which is an \(E\)-vector space. Newton identities now give
the relative coefficients \(a_j\) in \(E\). At each embedding they
are sums of products of weight-one Weil \(p\)-numbers, so
\(a_j\in\mathcal O_E\).

Pick an auxiliary rational prime \(\ell\) away from the order index
and the discriminant of \(E\). Put \(C_j=\binom hj^{2e}\). Choose \(N\)
with \(\ell^N>2\max_{1\le j<h}C_j\), and enlarge the base once to
make \(A[\ell^N]\) and the \(\ell^N\)-th roots of unity rational.
The Tate lattice is free of rank \(h\) over
\(\mathcal O_E\otimes\Z_\ell\) since each local factor is a discrete
valuation ring and the corresponding summand is torsion-free of that
rank. The Frobenius congruence on this lattice and the cyclotomic
congruence at good degree-one primes of residue characteristic
\(p\ne\ell\) give
\[
 a_j\equiv\binom hj\pmod{\ell^N(\mathcal O_E\otimes\Z_\ell)},
 \qquad p\equiv1\pmod{\ell^N}.
\]
On the complex side the Weil bound is
\[
 |\sigma(a_j)|\le\binom hj p^{j/2},
 \qquad |\Nm_{E/\Q}(a_j)|\le C_jp^{ej}.
\]
If \(p^{ej}\) divides the norm, then
\(p^{-ej}\Nm_{E/\Q}(a_j)\) is an integer in \([-C_j,C_j]\) congruent
to \(C_j\) modulo \(\ell^N\). Thus it equals \(C_j\). This is the
bounded congruence argument of Ogus \cite[Proposition~2.7.2]{ogus}
applied to the normalized norm.

Equality in the product bound implies equality at every embedding.
In a component, all \(\binom hj\) products of \(j\) eigenvalues have
the same modulus. Equality in their triangle inequality forces them
to agree. Given two distinct indices, we can choose \(j-1\) other
indices and compare the two products to see that the two eigenvalues
agree. This holds for all indices because \(1\le j<h\). The repeated
eigenvalue property carries to all realizations of each component.
At a split place of a mixed component the common eigenvalue is its
trace divided by \(h\) and is therefore in \(\Q_p\). Its valuation
must be \(d_\sigma/h\in(0,1)\), which is impossible in \(\Q_p\).
Consequently every intermediate norm has valuation at most \(ej-1\).
In particular all these coefficients are nonzero.

In a split component the least possible valuation of a \(j\)-fold
eigenvalue product is \(\sum_{k=1}^j s_{\sigma,k}\).
The valuation of its sum with the other such products is no smaller,
and is integral because the relative coefficient lies in \(\Q_p\).
Summing the rounded bounds proves the slope inequality. Absolute
degree-one primes of \(L\) have natural density one by the prime ideal
theorem, since the higher-degree primes of norm at most \(x\) lie
over rational primes at most \(\sqrt x\). This proves
Theorem~\ref{thm:bounds}.

\section{Sharp numerical bounds from the signature}\label{sec:numerical}

A numerical component of height \(h\) and dimension \(m\) is
a convex polygon with slopes in \([0,1]\), endpoints \((0,0)\)
and \((h,m)\), and integral breakpoints. Equivalently, each
reduced slope \(a/c\) occurs with multiplicity divisible by \(c\)
\cite[\S1.1]{oort}. The polygon order used here is pointwise order
on the prefix sums \(S_j\) of the ordered slopes. Duality sends
\(S_j\) to \(j-m+S_{h-j}\), so it preserves this order.

Orient the Hodge multiplicities as \((m_i,h-m_i)\) with
\(m_i\le h/2\), and let \(b\) be the number of positive \(m_i\).
The ordinary slack at the \(j\)-th norm bound, before the strict
unit margin, is
\begin{equation}\label{eq:ordinary-slack}
 D_j=\sum_i\min(j,h-j,m_i).
\end{equation}
Indeed the ordinary two blocks have a total rounded prefix of
\(j-\min(j,h-j,m_i)\). If \(b>0\), an excess \(\Delta_j\) above
the ordinary rounded prefix is allowed precisely when
\(\Delta_j\le D_j-1\). The all-definite case is ordinary immediately.

\begin{proof}[Proof of Corollary~\ref{cor:ranktwo}]
The sum of the two slopes on each balanced side is one.
An ordinary balanced pair has zero excess at the only index
\(j=1\). For a nonordinary pair both slopes on both sides are
strictly positive, so each rounded first prefix is one and the
excess is two. As the ordinary slack is \(b\), the strict margin
gives \(2r\le b-1\). Each such pair removes two units of
\(p\)-rank, giving the stated bound. With no balanced pair,
each split component is definite and ordinary.
\end{proof}

\begin{proof}[Proof of Corollary~\ref{cor:rankfour}]
In a nonordinary pair of type \((2,2)\), each side has fewer than
two zero slopes, since two zero slopes would force the other
two to be one. The sum of the two rounded prefixes at \(j=2\)
is therefore at least two on this pair, compared with zero in
the ordinary case. Its slack is two, and definite pairs add
no slack. This violates the strict unit margin.
\end{proof}

\begin{theorem}[Least polygon with fixed endpoint multiplicities]\label{thm:extremal}
Fix \(h,m\), and fix the numbers \(z,o\) of zero and unit slopes
of a numerical component. Set \(q=h-z-o\), \(w=m-o\).
If \(q=0\) the polygon is ordinary. Otherwise \(0<w<q\).
The least possible interior polygon is the chord from \((0,0)\) to
\((q,w)\) when \(w=1\) or \(q-w=1\). In the remaining cases it
joins \((0,0)\), \((q-w,1)\), and \((q,w)\), with interior slopes
\[
 \left(\frac{1}{q-w}\right)^{q-w},\qquad
 \left(\frac{w-1}{w}\right)^w .
\]
Consequently the least possible paired excess at each \(j\)
is attained simultaneously and is
\begin{equation}\label{eq:interval}
  \mathbf 1_{z<j<h-o}+\mathbf 1_{o<j<h-z}.
\end{equation}
The same expression is zero for ordinary endpoints.
\end{theorem}

\begin{proof}
Write \(a=q-w\) and \(c=w\). Each proper interior vertex
\((x,y)\) has integer coordinates with \(y\ge1\) and
\(y\ge x-a+1\). The second bound follows because the remaining
slopes are strictly less than one. If \(c=1\) there is no such
vertex because its ordinate would be strictly between zero and one.
If \(a=1\) the same holds after dualizing.

For \(a,c\ge2\), the displayed two slopes are increasing, with
equality possible when \(a=c=2\), in which case the middle point is
collinear. The resulting lower convex polygon is the maximum of
the lines \(x/a\) and \(x-a+1-(x-a)/c\).
At any proper vertex the first line is bounded by \(1\) if
\(x\le a\), and by \(x-a+1\) if \(x\ge a\).
The second line is bounded by \(1\) when \(x\le a\), and by
\(x-a+1\) otherwise. Thus both lie below every vertex. Convexity
puts the proposed polygon below every joining segment, including
the isoclinic chord.

Restore the zero and unit tails. From \(j=z+1\) through
\(j=h-m=z+a\), the first interior segment has rounded prefix one.
At \(j=h-m+t\), \(1\le t<c\), the next segment has rounded prefix
\(1+t\), while at \(t=c\) the prefix equals \(c\). Hence excess
over the ordinary prefix is one precisely for \(z<j<h-o\).
If \(a=1\) or \(c=1\), taking the ceiling of the chord gives
the same strict interval. Dualizing swaps \(z,o\). Since duality
preserves order, the minimizer supplies both terms in
\eqref{eq:interval} simultaneously.
\end{proof}

The symmetric specialization of this polygon, with equal numbers
of extreme slopes and equal residual dimension and codimension,
appears in Achter and Pries \cite[Definition~2.4]{achterpries}.
The oriented component statement gives an endpoint-only feasibility
test for the strict simultaneous bounds. For each pair select
either the ordinary endpoints or integers
\(0\le z\le h-m-1\), \(0\le o\le m-1\).
Test the sum of \eqref{eq:interval} against \(D_j-1\) and minimize
\(\sum_i(z_i+o_i)\). Any feasible endpoints lift to numerical
polygons by Theorem~\ref{thm:extremal}, and every admissible
collection gives feasible endpoints. This proves exactness.

\begin{proposition}[When the inequalities force ordinary reduction]\label{prop:signatures}
The only signatures for which the numerical restrictions force the
ordinary polygon are all-definite signatures, one mixed pair of type
\((1,h-1)\), two mixed pairs in rank two, and one mixed pair of
type \((2,2)\) in rank four. Definite pairs can be appended freely
to all cases.
\end{proposition}

\begin{proof}
If a lone mixed pair has \(m=1\), then at \(j=1\) the slack
is one. A nonordinary block on its dimension-one side cannot
have a unit slope, so the complementary block has no zero, and
the paired rounded prefix increases by at least one.
The rank-two and rank-four exceptions follow from the arguments
for Corollaries~\ref{cor:ranktwo} and~\ref{cor:rankfour}.

For the converse, in one mixed pair replace the last zero and
first unit slope by two copies of \(1/2\), preserving duality.
This is the least nonordinary paired polygon for that pair.
In fact a nonordinary block with endpoint losses \(a,c\ge1\) has
prefix at the ordinary kink \(h-m\) at least \(1/2\) by
Theorem~\ref{thm:extremal}, while its other integer prefixes
are at least ordinary. The replacement has excess
\[
        \Delta_j=\mathbf 1_{j=m}+\mathbf 1_{j=h-m}
\]
by \eqref{eq:interval}. For \(h=2\) and \(b\ge3\), the slack is
at least three at its only index, so the replacement is allowed.
For \(h>2\) and at least two mixed pairs, choose any mixed pair.
If its two affected indices are distinct, the slack is at least
two at each. If they coincide, then \(m=h/2\ge2\), and at that
index the modified pair contributes at least two to the slack
while a second pair contributes at least one. Finally if there is
one mixed pair with \(m\ge2\), the slack at both affected indices
is \(m\). The replacement fails the strict unit margin precisely
when they coincide at \(m=2\), the already listed rank-four case.
This gives a witness in every other signature.
\end{proof}

The total number of signatures up to dimension \(G\) is
\[
 \sum_{h=1}^G
 \left(\binom{\lfloor h/2\rfloor+\lfloor G/h\rfloor+1}
                    {\lfloor h/2\rfloor+1}-1\right)
\]
by counting multisets of \(e\) multiplicities and summing over \(e\).
For \(G=30\) this is 1188. The four disjoint cases of
Proposition~\ref{prop:signatures} give 111, 81, 14 and 7,
respectively. The first count is
\(\sum_{h=1}^{30}\lfloor30/h\rfloor\), the second omits \(h=1\),
and the remaining cases count \(2\le e\le15\) and \(1\le e\le7\).

\begin{theorem}[Sharp numerical \(p\)-rank]\label{thm:sharp-prank}
If \(b=0\), the minimum is \(eh\). Otherwise put
\[
 H=\lfloor h/2\rfloor,\qquad L_j=\max(0,2b+1-D_j)
 \quad (1\le j\le H).
\]
The sharp minimum is \(h(e-b)+M\), where the mixed contribution
\(M\) is as follows.
\begin{itemize}
\item If all mixed multiplicities are one, then
\[
 M=\begin{cases}
  2\lceil(b+1)/2\rceil, & h=2,\\
  h+H(b-1), & h\text{ odd},\\
  h+(H-1)(b-1)+\lceil(b-1)/2\rceil, & h\ge4\text{ even}.
 \end{cases}
\]
\item If \(h=4\) and all mixed multiplicities are two, then \(M=b+3\).
\item In every other case with some mixed multiplicity at least two,
\[
 M=\begin{cases}
  \displaystyle\sum_{j=1}^H L_j, & h\text{ odd},\\
  \displaystyle\sum_{j=1}^{H-1}L_j+\lceil L_H/2\rceil,
      & h\text{ even}.
 \end{cases}
\]
\end{itemize}
\end{theorem}

\begin{proof}
For a mixed pair write \(P_j=2-\Delta_j\), using the minimizing
endpoints. The simultaneous condition is
\(\sum P_j\ge L_j\) for \(1\le j\le H\).
For a nonordinary pair, \(o<H\) and
\[
 P_j=\mathbf 1_{j\le z}+\mathbf 1_{j\le o}
             +\mathbf 1_{j\ge h-z}.
\]
For an ordinary pair, \(P_j=2\).
If \(h\) is odd, summing \(P_j\) gives exactly \(z+o\) for a
nonordinary pair and \(h-1\) for an ordinary pair. If \(h\) is even,
give the index \(H\) weight one half. This weighted sum has the
same two values, and \(P_H\) is always even.
These observations give the lower bounds in the last case.

For sharpness there, use only nonordinary pairs. Think of the two
endpoint counts as slots of capacities
\(\min(H,h-m_i-1)\) and \(m_i-1\).
Choose lengths in the slots with prescribed tail counts \(L_j\)
for \(j<H\), and \(L_H\) in odd rank or \(\lceil L_H/2\rceil\)
in even rank at \(H\). At these lengths the third indicator in
\(P_j\) is absent except at the even midpoint, where it doubles
the contribution of a slot of length \(H\).
For a decreasing list of tail counts a feasible assignment
exists if each count is at most the number of slots with capacity
at least its index. To see this, sort the lengths and capacities
and match in decreasing order.

At \(j=1\) there are \(b\) zero slots and at least one unit slot.
For \(2\le j<H\), and for \(j=H\) in odd rank, the \(b\)
zero slots suffice since \(D_j\ge b+1\).
At the even midpoint only unbalanced pairs supply a zero slot.
Let their number be \(u\). For \(H\ge3\), if there is any
balanced pair, then \(D_H\ge u+3(b-u)\) and \(L_H\le u\).
If there is none, \(u=b\) already suffices. For \(H=2\) in
this case \(u\ge1\), \(L_H=u+1\), and
\(\lceil(u+1)/2\rceil\le u\). The resulting assignment
attains the lower bounds.

If all mixed multiplicities are one and \(h\ge3\), at least
one pair is ordinary. Otherwise all \(o=0\) and each \(P_1\)
is at most one, whereas \(L_1=b+1\).
With \(c\ge1\) ordinary pairs, each index needs at least
\(T=\max(0,b+1-2c)\) further units of \(P_j\).
In odd rank choose \(T\) remaining zero endpoints of length \(H\).
In even rank use \(\lceil T/2\rceil\) of length \(H\) and the
remaining \(T-\lceil T/2\rceil\) of length \(H-1\).
There are enough slots because \(T\le b-c\), and they attain
the weighted lower bound for this \(c\). Increasing \(c\) costs
\(h\) and can save at most \(h-1\) of this lower bound, so
\(c=1\) is optimal. Rank two follows directly by counting
ordinary pairs.

Finally in rank four with all mixed multiplicities two,
every nonordinary pair has \(P_2=0\), so at least one
ordinary pair is necessary. With \(c\) such pairs,
the lower bound is \(4c+\max(0,b+1-2c)\), from index one.
It is minimized at \(c=1\), and is attained with one endpoint
of length one in each of the other pairs.
\end{proof}

The dimension-30 table includes exact polygon multiplicities
as well as endpoint minima. The endpoint formula gives feasibility
and optimization. Individual polygons with fixed endpoints can
differ in admissibility. For instance, at \(h=8,m=4,z=2,o=1\)
the extremal interior \((1/2)^2,(2/3)^3\) passes with one mixed
pair, while the alternative \((3/5)^5\) fails at \(j=4\).
For full multiplicities one can transfer the paired rounded
profiles of integral convex paths instead. The verification
replays retain both the numerical minima and the historical full
counts. The closed formula gives minima for \(5{,}475{,}208\)
signatures through dimension 200, of which 2145 force ordinarity.
The companion verification archive contains the dimension-30 table,
the larger summary, and the replays for both computations.


\section{The first case with three balanced pairs}\label{sec:residual}

In rank two write \(\tau=a_1\), \(\delta=a_2\).
Polarized Frobenius gives \(\bar\tau=p\tau/\delta\). At a split
place of dimension \(d_\sigma\), one has \(v_p(\delta)=d_\sigma\).
It follows that for a pair the sum of the two trace valuations is
\[
 \big|1-d_\sigma\big|
   +2\min\{v_{\sigma}(\tau),v_{\bar\sigma}(\tau)\}.
\]
With exactly three balanced pairs, a single nonordinary pair forces
\(v_p(\Nm(\tau))=e-1\). Indeed it supplies at least two valuations
in addition to the \(e-3\) from definite pairs, and
Theorem~\ref{thm:bounds} allows at most \(e-1\). Both bad traces
have valuation exactly one. The remaining balanced traces and
the dimension-zero traces in definite pairs are units.

For the six-dimensional signature in which all three pairs are
balanced, this pattern also identifies a limit of related trace
tests. Give \(X=\tau\) weight one and \(Y=\delta\) weight two.
Fix a nonzero homogeneous polynomial
\[
             P(X,Y)=\sum_k c_kX^{w-2k}Y^k\in E[X,Y]
\]
of weight \(w\), and let \(k_0,k_1\) be the least and greatest
indices with \(c_k\ne0\). At all but finitely many rational primes
its nonzero coefficients are units. At each ordinary component the
unique least term has valuation \(k_0\). At each nonordinary
component the unique least term has valuation \(w-k_1\).
Consequently the norm has valuation
\[
       4k_0+2(w-k_1)\le 2w+2k_0<3w
\]
if \(P\) has any trace term. The final inequality uses
\(2k_0<w\). The exponent \(3w\) is the normalization suggested by
its Weil weight. In particular a fixed homogeneous polynomial in
the two relative coefficients, including the trace of a positive
symmetric power, adds no divisibility contradiction for this
polygon. The determinant alone has integral normalization in both
the ordinary and nonordinary cases. This argument leaves open
whether the surviving pattern has positive density for any such
fixed abelian variety.

\subsection{Realization over prime fields}

The rank-two bound is sharp for prime-field data even when the
auxiliary torsion level is fixed arbitrarily high.

\begin{proposition}[Finite-field sharpness]\label{prop:sharpness}
Fix a totally real field \(F\) of degree \(e\), put \(E=F(i)\),
and fix an odd rational prime \(\ell\) and \(N\ge1\). For each
\(0\le r<e/2\), and for every sufficiently large rational prime
\(p\) split completely in \(E\) and congruent to \(1\) modulo
\(\ell^{2N}\), there is an abelian variety over \(\mathbb F_p\)
of dimension \(2e\) with the following properties. It admits a
Rosati-compatible rank-two action by \(E\), all \(e\) split pairs
have dimensions \((1,1)\), exactly \(r\) pairs are nonordinary,
and its \(\ell^N\)-torsion is rational.
\end{proposition}

\begin{proof}
Let \(M=\ell^{2N}\) and let \(I\subset\mathcal O_F\) be the product
of any \(r\) distinct primes over \(p\). We first choose
\(\tau\in I\) with \(\tau\equiv 2\pmod{M\mathcal O_F}\) and
\(\Q(\tau)=F\), such that \(|\sigma(\tau)|\ll_{F,M}p^{r/e}\)
for each real embedding \(\sigma\).

Here is the lattice argument. Write \(I=\beta J\), where \(J\) comes
from a fixed finite list of fractional ideal class representatives.
The logarithmic unit lattice of \(F\) has a bounded fundamental
domain in the hyperplane of coordinate sum zero. Multiplying
\(\beta\) by a unit lets us bound each \(|\sigma(\beta)|\) by
a constant times \(|\Nm_{F/\Q}(\beta)|^{1/e}\), hence by a constant
times \(p^{r/e}\). Transport a fixed \(\Z\)-basis of the lattice
\(J\) to a basis of \(I\). All its coordinates in the real embeddings
are bounded by \(O_F(p^{r/e})\). Since \(I\) is coprime to \(M\), the
congruence is a nonempty coset modulo \(MI\) in \(I\). Reducing a
representative in a fundamental parallelepiped gives the required
size bound. If \(F\) has \(s\) proper subfields, vary this
representative by \(MI\)-basis combinations with coefficients in
\(\{0,\ldots,s\}\). In grid coordinates each proper subfield
meets at most \((s+1)^{e-1}\) points, since it lies in a
proper affine hyperplane. Some point avoids all the subfields.
This preserves the bound with a larger fixed constant.

For sufficiently large \(p\), the strict inequality
\(|\sigma(\tau)|<2\sqrt p\) follows from \(r/e<1/2\).
The norm is nonzero and bounded in absolute value by
\(O_{F,M}(p^r)\), so \(p^{r+1}\) cannot divide it. Thus
\(v_{\mathfrak q}(\tau)=1\) at the selected primes, and
\(\tau\) is a unit at the other primes of \(F\) over \(p\).
Let \(\pi\) be a root of \(X^2-\tau X+p\). This is a Weil
\(p\)-integer, since the quadratics in all real embeddings
have nonreal conjugate roots of modulus \(\sqrt p\).
The field \(C=\Q(\pi)=F(\pi)\) is CM of degree \(2e\),
using \(\tau=\pi+p/\pi\). At the selected primes the
quadratic is Eisenstein with slopes \((1/2,1/2)\).
At every other prime over \(p\) it splits by Hensel's lemma
with slopes \((0,1)\).

By Honda--Tate \cite[Theorem~1]{tate}, \(\pi\) is the Frobenius
of a simple abelian variety \(B/\mathbb F_p\) up to isogeny.
The only potentially fractional local endomorphism invariant
over \(p\) has the form \(2\cdot(1/2)=1\) modulo \(1\).
All invariants vanish, so the rational endomorphism algebra over
\(\mathbb F_p\) is \(C\), and \(B\) has dimension \(e\).
By Waterhouse \cite[Theorem~6.1]{waterhouse}, choose \(B\) with
endomorphism ring \(\mathcal O_C\).
The number \(\eta=(\pi-1)/\ell^N\) belongs to that ring, since
it satisfies the monic equation
\[
 \eta^2+\frac{2-\tau}{\ell^N}\eta
       +\frac{1-\tau+p}{\ell^{2N}}=0
\]
over \(\mathcal O_F\). Therefore Frobenius acts trivially on
\(B[\ell^N]\).

On \(A=B^2\), let \(F\) act diagonally and let \(i\) act by
the matrix with rows \((0,-1)\) and \((1,0)\).
On \(B\), Rosati sends \(\pi\) to \(p/\pi\), so it fixes
\(F=\Q(\pi+p/\pi)\). The product Rosati involution fixes the
diagonal action and sends the rotation matrix to its negative.
Over \(\Q_p\), the height-two \(F\)-components of \(B\)
follow from the relative quadratic Frobenius polynomial.
The action of \(i\) splits at both places above each prime of \(F\).
Each of the two eigenspaces inherits the height-two component
of \(B\) at that prime, with slope sum one. Thus exactly
the selected \(r\) pairs are nonordinary, and all \(\ell^N\)-torsion
is rational as claimed.
\end{proof}

The primes in the proposition exist by Chebotarev. Taking
\(r=\lfloor(e-1)/2\rfloor\) realizes the bound for every all-balanced
signature over prime fields. The examples vary with \(p\).
A conclusion about the density of the residual pattern for a fixed
characteristic-zero variety requires global input beyond these
split Frobenius and torsion data.

\section{A global criterion for three balanced pairs}\label{sec:linked}

We now keep the abelian variety fixed and take relative rank two
with exactly three balanced pairs. After extending the base,
assume it contains a normal closure of \(E\), all the geometric
endomorphisms are defined, and the \(\ell\)-adic algebraic monodromy
group \(G\) is connected. Its rank-two summands under the embeddings
of \(E\) are denoted by \(V_\sigma\).
The regular class function
\[
 B_\sigma(g)=\frac{\operatorname{tr}(g\mid V_\sigma)^2}
                  {\det(g\mid V_\sigma)}
\]
is the character of the adjoint representation plus one. Call
two embeddings linked when these functions agree on \(G\).

\begin{theorem}[Linked projective monodromy]\label{thm:linked}
If two embeddings in different conjugate pairs are linked, then
\(A\) has ordinary reduction at a density-one set of primes
after finite extension.
\end{theorem}

\begin{proof}
At good degree-one primes the relative trace \(\tau\) and determinant
\(\delta\) are algebraic and independent of cohomology, by the
comparison in Section~\ref{sec:coefficients}. Thus equality of two
class functions is equivalent to equality of the corresponding
conjugates of \(\tau^2/\delta\) at all such primes. The reverse
implication follows from Chebotarev: the complement of a class
function's zero set in the Galois image, if nonempty, contains an
open conjugacy-invariant set and is detected in a finite quotient.
The absolute degree-one primes have density one. In particular
linkage is invariant under the rational Galois group on the
embeddings. Polarization gives
\(\bar\tau=p\tau/\delta\) and \(\bar\delta=p^2/\delta\),
so conjugate embeddings are always linked.
On the conjugate pairs, every equivalence class therefore has the
same cardinality, by Galois transitivity.

Take the degree-one primes of the extension in
Theorem~\ref{thm:bounds}, enlarged as above. The trace is nonzero
there. At each split place, the valuations of the projective
trace are \(2v_p(\tau_\sigma)-d_\sigma\); they are odd on
balanced pairs and even on definite pairs. Thus every class is
entirely balanced or entirely definite. A nontrivial linkage
makes the common cardinality greater than one. It divides three
by the number of balanced pairs and hence equals three.
At a balanced component the trace is a unit precisely when the
slopes are \((0,1)\), so equality of the three projective trace
valuations makes all failures simultaneous. By
Corollary~\ref{cor:ranktwo} at most one balanced pair can fail.
Consequently all the selected primes are ordinary.
\end{proof}

Equivalently, the field generated by the good degree-one values
\(\tau^2/\delta\) lies in \(E^+\), and the linkage classes of
conjugate pairs are the fibres of restriction to that field.
Thus a proper projective trace field in \(E^+\) already
implies the conclusion.

\begin{proposition}[Residual geometric endomorphism types]\label{prop:end-frontier}
In the case of three balanced pairs, exclude the potentially CM
varieties and the varieties satisfying Theorem~\ref{thm:linked}.
Up to geometric isogeny the remaining possibilities are:
\begin{enumerate}
\item an absolutely simple variety with full endomorphism algebra
      \(E\) and \(e\ge4\);
\item the square of an absolutely simple threefold whose geometric
      endomorphism field is totally real cubic;
\item an absolutely simple sixfold with a totally indefinite
      quaternion division endomorphism algebra over a totally
      real cubic field.
\end{enumerate}
In these cases the derived monodromy projections on the conjugate
pairs are independent copies of \(\mathrm{SL}_2\), up to central
isogeny.
\end{proposition}

\begin{proof}
Each geometric isotypic component inherits a faithful \(E\)-action.
If there is more than one, the total relative rank two splits as
one plus one, which gives potential CM. Thus write \(A\sim C^n\)
with \(C\) absolutely simple, and let \(F\) and \(t\) be the center
and Schur index of its geometric division endomorphism algebra.

By Faltings's semisimplicity and endomorphism comparison
\cite[Theorems~3--4]{faltings}, over
\(\overline{\Q}_\ell\) the monodromy commutant has matrix blocks of
size \(nt\) indexed by embeddings of \(F\). The corresponding
irreducible monodromy constituents all have the same dimension.
Indeed the \(F\)-component dimension is uniform by the rational
Betti structure. The common constituent dimension is one or two,
as the \(E\)-summands have dimension two. If it is one, the
rational endomorphism algebra admits a commutative semisimple
subalgebra of degree \(2\dim A\) by taking a maximal commutative
\'{e}tale algebra in each matrix factor. This is the potential-CM
case. Such varieties are ordinary at degree-one primes after
adjoining the normal closures of their CM fields of maximal
degree, since each corresponding component is height one.

Otherwise each \(V_\sigma\) is irreducible. The centralizer of
\(E\) in the geometric endomorphism algebra is exactly \(E\),
as its split centralizer has dimension one per \(V_\sigma\).
Therefore \(F\subset E\) and \([E:F]=nt\).
All embeddings above an embedding of \(F\) give isomorphic
representations, by the above commutant decomposition. In the
absence of a cross-pair link, a fibre has size at most two,
and a fibre of size two consists of conjugates.
If the size is two, \(F=E^+\). Equal determinant valuations
at the two embeddings above each place then force every
conjugate pair to be balanced. Hence \(e=3\).
Here \(nt=2\), giving either the real-cubic square or a quaternion
division endomorphism algebra. If the latter is ramified at a
real place, positivity forces the Rosati involution there to be
canonical quaternion conjugation, of symplectic type. An
orthogonal involution would fix a nonzero pure quaternion and
therefore violate positivity. After scalar extension the relevant
polarization on the two-dimensional matrix module and its
two-dimensional multiplicity space is the tensor of an alternating
and a symmetric form. The derived monodromy would then act through
\(\mathrm{SO}_2\) on the multiplicity space, contradicting its
irreducibility since the connected center acts by scalars.
Thus the quaternion algebra is totally indefinite.

A fibre of size one gives \(n=t=1\) and full endomorphism
algebra \(E\). If \(e=3\), all pairs would be balanced and the
relative determinant would have type \((1,1)\) at every embedding.
It is therefore trivial on the Hodge circle, and its kernel is
defined over \(\Q\), so it is trivial on the Hodge group.
Polarization and the self-duality of a two-dimensional representation
of determinant one make the conjugate summands isomorphic as Hodge
group modules. The complex Hodge commutant then contains
off-diagonal maps, whereas the commutant of the rational Hodge
structure equals the algebraic endomorphism algebra of the
abelian variety after tensoring with \(\Q\). This excludes \(e=3\)
in the full-field case.

For any of the remaining types, each derived projection is
\(\mathrm{SL}_2\) since the monodromy is connected reductive and
the two-dimensional summand is irreducible. A nontrivial
projection of the semisimple derived Lie algebra onto
\(\mathfrak{sl}_2\) uses one simple ideal. Two projections
using the same ideal give the same projective trace function,
as they are both the standard two-dimensional action, up to
conjugacy, and scalars cancel. The absence of a cross-pair link
therefore gives distinct factors for distinct pairs.
\end{proof}

In particular, a larger endomorphism algebra with CM center
forces ordinarity in this setting. The remaining real-cubic
square case already contains the ordinary-reduction question for
real-multiplication threefolds over arbitrary number fields.
For such a polarized threefold \(C\) with real field \(F\) fixed
by Rosati, the product \(C^2\) carries the diagonal \(F\)-action
and the standard rotation giving \(F(i)\). The three real
embeddings of \(F\) have equal complementary Hodge dimensions,
so this action has exactly three balanced pairs. Ordinarity of
\(C^2\) and of \(C\) are equivalent. Fité's cubic result
\cite[Theorem~3]{fite} assumes the real field is not contained
in a field of definition of its action, and the general theorem of
Wang and Zhang \cite[Theorem~1.3]{wangzhang} assumes an
auxiliary splitting pattern. These hypotheses need not hold
in the split frontier above.


\section{Comparison and a cover example}

Cantoral Farf\'an, Li, Mantovan, Pries and Tang use norms of selected
relative exterior coefficients in their simple-signature criterion
\cite[Assumptions~2.1 and~2.5, Proposition~3.12 and Theorem~4.18]{cantoral}.
Their eventual ordinary theorem covers one mixed pair of type
\((1,h-1)\) when the geometric endomorphism algebra equals the acting
CM field. Theorem~\ref{thm:bounds} allows an embedded CM field and
arbitrary signatures; it imposes simultaneous constraints even where
they do not force ordinary reduction. Their conclusions about
\(\mu\)-ordinarity and about density over the original field in
specific cases use different hypotheses.

Fit\'e bounds the norm of the relative trace for an embedded Galois
field of degree \(g\) and obtains a positive-density
\(p\)-rank bound \(\lceil g/2\rceil\)
\cite[Proposition~6 and Theorem~2]{fite}. The proof obtains density
one over a suitable extension. With a rank-two Rosati-compatible
CM action, its weak valuation bound gives
\(r\le\lfloor b/2\rfloor\). Our strict inequality saves a pair
when \(b>0\) is even and applies without the Galois assumption.
Fit\'e's fourfold with imaginary quadratic signature \((2,2)\)
already gives the case \(e=1\) of Corollary~\ref{cor:rankfour}
\cite[Theorem~4 and \S8]{fite}. His exterior-square argument
uses weight-two saturation and the resulting scalar-eigenvalue
condition. Sawin uses weight-zero twists of an exterior-square
trace for abelian surfaces \cite{sawin}. Wang and Zhang use
normalized relative Frobenius norms for GL\(_2\)-type varieties
under different base-field and splitting hypotheses. Their
general number-field theorem excludes Albert type IV
\cite[Theorems~1.1--1.3]{wangzhang}.

The rank-four signature occurs in every dimension \(4e\). For a CM
field of degree \(2e\), take two CM types differing at just one
conjugate pair. The product of two copies of each of their CM
abelian varieties, with diagonal field action and product
polarization, has the required signature. This example also has
additional endomorphisms and its ordinarity follows from CM theory.

\begin{example}
Consider the primitive conductor-eight factor of the Jacobian of
a cyclic cover of the line with six distinct branch points and
exponents \((1,1,1,1,2,2)\) modulo eight. For a character \(k\)
prime to eight, the coherent eigenspace dimension in one consistent
character convention is \(-1+\sum_i\{ka_i/8\}\), where
\(\{\cdot\}\) denotes fractional part
\cite[\S2]{moonen}. For \(k=1,3,5,7\) this gives
\((0,2,2,4)\), so the factor has dimension eight and a rank-four
\(\Q(\zeta_8)\) action with one \((2,2)\) pair. The deck action
preserves the polarization and gives the required Rosati involution.
Corollary~\ref{cor:rankfour} applies to every such cover over a
number field, including fibers with additional endomorphisms.
\end{example}

\section{An application to polynomial moments}

A normalizer of the acting CM field can strengthen the rank-two
bound. Suppose a group of invertible geometric quasi-endomorphisms normalizes
\(E\). After defining representatives of its finite image in
\(\operatorname{Aut}(E)\) over a finite base extension, the
embeddings in the same orbit have isomorphic Frobenius components
at good split primes outside finitely many exceptions.
Consequently the nonordinary balanced pairs form a union of orbits
on conjugate pairs. For \(b=3\), for example, a transitive action
on its three pairs forces ordinarity. More generally, simultaneous
vanishing of coherent Frobenius in \(\lceil b/2\rceil\) distinct
rank-two balanced pairs contradicts Corollary~\ref{cor:ranktwo}.
On either side of a balanced pair the coherent space has
dimension one, so its zero Frobenius scalar makes the pair
nonordinary. A single isolated moment coefficient usually supplies
just one such scalar.

The arbitrary-rank result also gives a moment obstruction without
a symmetry of the cover. Write \(f\) on the sphere quadric
\(t^2+2xy=1\) as a Laurent polynomial in \(x,t\), by substituting
\(y=(1-t^2)/(2x)\). Its sphere moment is
\[
 \mathcal A(f^m)=\operatorname{CT}_x\,\frac12
                       \int_{-1}^1 f(x,t)^m\,dt ,
\]
interpreted as a rational linear functional on the Laurent chart.
Assign weight \(2i+j\) to \(x^it^j\).

\begin{proposition}[A squarefree sextic face]\label{prop:sextic}
Suppose the greatest Laurent weight of \(f\) is four, with
initial form \(t^4z^{-2}P(z)\), where \(z=x/t^2\) and \(P\)
has degree four, is squarefree, and has nonzero constant term.
Then over any characteristic-zero field some positive power
has nonzero sphere moment. All smaller Laurent weights are
allowed. If the phase has normal-form degree at most six,
its least sphere degree is exactly six.
\end{proposition}

\begin{proof}
If the moments all vanish, specialize the finitely generated
coefficient algebra over \(\Q\) to a number field after inverting
the leading and constant coefficients of \(P\) and its discriminant.
The vanishing identities and the specified initial form survive.
At a good prime of degree one with rational characteristic
\(p=4k+1\), the largest possible \(t\)-degree in the angular
constant term of \(f^k\) is \(4k=p-1\).
Integration of an even \(t\)-degree \(j\) divides by \(j+1\)
and odd degrees integrate to zero. Multiplying by \(p\) and
reducing modulo the prime therefore gives
\[
       p\mathcal A(f^k)\equiv [z^{2k}]P(z)^k\pmod p .
\]
Only the initial form can contribute to this top degree.

The normalized cyclic cover \(W^4=z^2P(z)\) has branch
exponents \((2,1,1,1,1,2)\), including infinity.
Its primitive order-four Jacobian factor has \(\Q(i)\)-rank
four and coherent signature \((1,3)\), by the fractional-part
formula. It is ordinary at the degree-one primes supplied by
Proposition~\ref{prop:signatures} after finite extension.
Here the deck action again supplies the polarization-compatible
field, and the degree-one split primes satisfy \(p=4k+1\).
The coherent \(H^1\) comparison of the curve with its Jacobian
is equivariant for the deck action.
For the character of \(W\), the coherent eigensheaf has generators
\(W\) on the finite chart and \(z^{-2}W\) at infinity. The
\v{C}ech quotient has basis the class of \(z^{-1}W\).
Its Frobenius image is
\(z^{-p}W^p=z^{-p+2k}P(z)^kW\); the coefficient of the basis class
is exactly \([z^{2k}]P(z)^k\). Ordinary Frobenius is bijective
on coherent \(H^1\), so this scalar is nonzero at these primes.
This contradicts purity.

In the normal form, the negative angular exponents come from
monomials involving \(y\), with at most one of \(x,y\) present.
The \(x^{-2}t^8\) term in the initial form comes from \(y^2t^4\)
with nonzero coefficient. This term has degree six and cannot
be supplied by a term of lower degree. Thus the stated normal-form
degree bound gives least sphere degree six.
\end{proof}

\subsection{From Hasse--Witt restraints to endpoint bounds}

At a split prime the stable rank of coherent Frobenius on an
embedding component is its slope-zero multiplicity. In
particular, the rank of the first Frobenius map bounds this
multiplicity from above \cite[Definition~2.1 and Lemma~2.2]{bouw}.
Using the complementary component gives an upper bound on the
unit-slope multiplicity as well.

\begin{proposition}[Testing simultaneous rank restraints]\label{prop:rankcaps}
Suppose a split mixed pair has endpoint restraints \(z\le R\)
and \(o\le R'\), obtained from coherent ranks on the corresponding
components. If the ordinary endpoints fit both restraints, use
them. Otherwise set
\[
       z=\min(R,h-m-1),\qquad o=\min(R',m-1).
\]
For all pairs, the simultaneous inequalities are numerically
compatible with the given endpoint restraints exactly when
these choices pass \eqref{eq:interval} with the strict unit
margin. If a definite pair violates its endpoint restraints
there is no compatible polygon.
\end{proposition}

\begin{proof}
The ordinary endpoints give zero excess. Once they are
unavailable, a nonordinary choice is necessary, and the proposed
values are coordinatewise largest among allowable endpoints.
Each interval in \eqref{eq:interval} can only shrink as either
endpoint multiplicity increases. Thus if the displayed choices
fail the test, every collection of slopes fails it. If they pass,
Theorem~\ref{thm:extremal} realizes the simultaneous numerical
inequalities.
\end{proof}

For \(h\ge3\), a vanishing coherent line on a pair of type
\((1,h-1)\) can be matched numerically by slopes
\(0^{h-2},1/2,1/2\) on its dimension-one side. The paired excess
is two at no index, and the ordinary slack of this pair is one at
every index. Thus if another mixed pair stays ordinary, the
chosen pair can fail while all inequalities hold. A test of one
line typically needs another computation. The next examples
supply one explicitly.

\subsection{Arbitrarily large supporting faces}

The prime-isolation argument also works with a prime power.
If the highest angular-constant height at a chosen moment order
is \(p^r-1\), and the Laurent coefficients are integral at the
prime, multiplying the moment by \(p^r\) isolates this highest
coefficient modulo the prime. Indeed every smaller height has
interval denominator of valuation less than \(r\).
For the Kummer model \(W^n=F(z)\) whose selected coherent line
has generator \(W/z\), Frobenius to the \(r\)-th power
acts on this line when \(p^r-1=nk\). Its scalar is
\([z^{p^r-1}]F(z)^k\), by projecting
\(z^{-p^r}W^{p^r}=z^{-p^r}F(z)^k W\) onto the generator.

\begin{proposition}[Powering a sextic face]\label{prop:scaled}
For every \(L\ge1\), the highest-weight form
\[
       t^{4L} z^{-2L}P(z)^L,\qquad z=x/t^2 ,
\]
with \(P\) as in Proposition~\ref{prop:sextic}, forces a nonzero
power moment, with arbitrary strictly lower-weight terms.
For phases of normal-form degree at most \(6L\) the least sphere
degree is exactly \(6L\).
\end{proposition}

\begin{proof}
Use the same connected conductor-four cover as in
Proposition~\ref{prop:sextic}; its Hasse coefficient is nonzero at
the good degree-one primes after one extension. Further adjoining
\(\zeta_{4L}\) keeps the conclusion at all but finitely many such
primes. With \(p=4Lk+1\), the moment at order \(k\) isolates
\([z^{2Lk}]P^{Lk}\), the original scalar at conductor-four
exponent \(Lk\). The negative endpoint has leading monomial
\(y^{2L}t^{4L}\) up to a nonzero scalar. Specialization from an
arbitrary characteristic-zero coefficient field works as in
Proposition~\ref{prop:sextic}.
\end{proof}

\begin{example}[A selected line on a nonordinary split cover]
Take \(P=(z-1)^2(z-5)^3\) and \(W^9=z^5P(z)\) over
\(\mathbb F_{73}\). Its primitive relative rank is two, with
three balanced pairs. The scalars on characters \(2\) and \(4\)
are zero and \(22\), respectively, by coefficient extraction
with the branch exponents reduced modulo nine. Explicitly, for
character \(i\) let \((a_0,a_1,a_5)\) be the least residues of
\((5i,2i,3i)\) modulo nine. The scalar is
\[
       [z^{72-8a_0}](z-1)^{8a_1}(z-5)^{8a_5}\pmod {73}.
\]
Thus the pair
\((2,7)\) fails and the pair \((4,5)\) is ordinary. Its unit
trace ensures that the relative trace norm is nonzero.
Two bad pairs would give norm divisibility by \(73^4\),
incompatible with the Weil bound \(64\cdot73^3\). This makes
the pair \((1,8)\) ordinary, so the selected scalar on character
one is nonzero, while the primitive factor is nonordinary.

Let \(73\nmid L\), and take highest weight \(9L\) with form
\(t^{9L}z^{-4L}P(z)^L\). With coefficients integral at 73,
all strictly lower weights are allowed, and the top form need
only have the specified reduction. Choose \(r\) with
\(73^r\equiv1\pmod {9L}\), and take moment order
\((73^r-1)/(9L)\). Prime-power isolation extracts the
\(r\)-fold scalar on the original connected cover.
It is nonzero, so the moment is nonzero. With normal-form
degree at most \(13L\), the negative endpoint gives
least sphere degree \(13L\). The exact coefficient
replay also gives the selected scalar \(11\bmod73\).
\end{example}

\begin{example}[Increasing connected conductor]
Set \(N=(37^r-1)/4\) for \(r\ge2\) and
\(P=(z-1)^2(z-3)^3\). At highest weight \(N\) with initial
form \(t^N z^{-4}P(z)\), the fourth moment is nonzero for all
strictly lower terms integral at 37. In fact it gives
\([z^{16}]P^4\equiv32\pmod {37}\) after multiplication by
\(37^r\). The connected cover \(W^N=z^{N-4}P(z)\) has this
selected scalar for Frobenius to the \(r\)-th power.

The primitive factor is nonordinary at 37. For a primitive
character \(i\), put \(x=i/N\). Its coherent dimension is
\(1+\lfloor4x\rfloor-\lfloor2x\rfloor-\lfloor3x\rfloor\);
this is two for \(1/4<x<1/3\) and one for \(0<x<1/4\).
Since \(9\mid N\), the character \(i=N/3-1\) is primitive.
It has ratio in \((1/4,1/3)\), while the representative
\(37^{r-1}i\equiv N/3-37^{r-1}\pmod N\) has ratio in
\((0,1/4)\). The unequal dimensions in a Frobenius orbit
prevent coherent Frobenius on the factor from being invertible.
The selected line still has nonzero iterated scalar. The
negative endpoint gives least sphere degree \(N+4\) if the
normal-form degree is at most \(N+4\).
\end{example}


\paragraph{AI assistance and review.}
ChatGPT (OpenAI) assisted with mathematical arguments, literature checks,
exposition, proof review and verification code.

\begin{thebibliography}{99}
\bibitem{achterpries}
J. Achter and R. Pries, Generic Newton polygons for curves of given
\(p\)-rank, in \emph{Algebraic Curves and Finite Fields},
Radon Ser. Comput. Appl. Math. \textbf{16} (2014), 1--21.
arXiv:\href{https://arxiv.org/abs/1311.5846}{1311.5846}.

\bibitem{bouw}
I. I. Bouw, The \(p\)-rank of ramified covers of curves,
\emph{Compos. Math.} \textbf{126} (2001), 295--322.
\href{https://doi.org/10.1023/A:1017513122376}{Sections 2 and 4--5}.

\bibitem{cantoral}
V. Cantoral Farf\'an, W. Li, E. Mantovan, R. Pries and Y. Tang,
\emph{Positive density of primes of ordinary reduction for abelian
varieties of simple signature},
arXiv:\href{https://arxiv.org/abs/2508.11174v1}{2508.11174v1}.

\bibitem{fite}
F. Fit\'e, Ordinary primes for some varieties with extra endomorphisms,
\emph{Publ. Mat.} \textbf{68} (2024), 27--40,
arXiv:\href{https://arxiv.org/abs/2005.10185v3}{2005.10185v3}.

\bibitem{faltings}
G. Faltings, Endlichkeitss\"atze f\"ur abelsche Variet\"aten
\"uber Zahlk\"orpern, \emph{Invent. Math.} \textbf{73} (1983), 349--366.
\href{https://doi.org/10.1007/BF01388432}{S\"atze 3--4}.

\bibitem{moonen}
B. Moonen, Special subvarieties arising from families of cyclic covers
of the projective line, \emph{Doc. Math.} \textbf{15} (2010), 793--819,
arXiv:\href{https://arxiv.org/abs/1006.3400}{1006.3400}.

\bibitem{oort}
F. Oort, Newton polygons and \(p\)-divisible groups: a conjecture by Grothendieck,
\emph{Ast\'erisque} \textbf{298} (2005), 255--269.
\href{https://www.numdam.org/item/AST_2005__298__255_0/}{\S1}.

\bibitem{ogus}
A. Ogus, Hodge cycles and crystalline cohomology, in
\emph{Hodge Cycles, Motives, and Shimura Varieties}, Lecture Notes in
Math. \textbf{900}, Springer, 1982, 357--414.
\href{https://publications.ias.edu/sites/default/files/Number45.pdf#page=376}{Proposition~2.7.2}.

\bibitem{pink}
R. Pink, \(\ell\)-adic algebraic monodromy groups, cocharacters,
and the Mumford--Tate conjecture,
\emph{J. Reine Angew. Math.} \textbf{495} (1998), 187--237.
\href{https://people.math.ethz.ch/~pinkri/ftp/AMT-v3.pdf}{Author copy}.

\bibitem{sawin}
W. F. Sawin, Ordinary primes for abelian surfaces,
\emph{C. R. Math. Acad. Sci. Paris} \textbf{354} (2016), 566--568,
arXiv:\href{https://arxiv.org/abs/1506.04784v2}{1506.04784v2}.

\bibitem{tate}
J. Tate, Classes d'isog\'enie des vari\'et\'es ab\'eliennes sur un
corps fini, S\'eminaire Bourbaki, Exp.~352 (1968--69), 95--110.
\href{https://numdam.org/item/SB_1968-1969__11__95_0/}{Theorem~1}.

\bibitem{wangzhang}
T. Wang and P. Zhang, Ordinary primes for GL\(_2\)-type abelian
varieties and weight 2 modular forms,
\emph{Mathematika} \textbf{72} (2026), e70081,
arXiv:\href{https://arxiv.org/abs/2503.21111v2}{2503.21111v2}.

\bibitem{waterhouse}
W. C. Waterhouse, Abelian varieties over finite fields,
\emph{Ann. Sci. \'Ec. Norm. Sup\'er.} \textbf{2} (1969), 521--560.
\href{https://www.numdam.org/item/ASENS_1969_4_2_4_521_0/}{Theorem~6.1}.
\end{thebibliography}
\end{document}
